\subsubsection{Distance point-segment / point-line}
\begin{lstlisting}[style=mystyle, language=C++]
// TC: O(1)
// usa: distancia minima de un punto a un segmento o a la recta que lo contiene
#include <bits/stdc++.h>
using namespace std;

// returns squared distance from p to line (a, b)
long long dist2PointLine(const Point& p, const Point& a, const Point& b){
	Point ab = b - a;
	long long area2 = abs((long long)ab * (p - a));
	long long len2 = 1LL * ab.x * ab.x + 1LL * ab.y * ab.y;
	if(len2 == 0) return 1LL * (p.x - a.x) * (p.x - a.x) + 1LL * (p.y - a.y) * (p.y - a.y);
	return (area2 * area2) / len2;
}

long long dist2PointSegment(const Point& p, const Point& a, const Point& b){
	Point ap = p - a;
	Point ab = b - a;
	long long dot1 = 1LL * ab.x * ap.x + 1LL * ab.y * ap.y;
	if(dot1 <= 0) return 1LL * ap.x * ap.x + 1LL * ap.y * ap.y;
	Point bp = p - b;
	long long dot2 = -1LL * ab.x * bp.x - 1LL * ab.y * bp.y;
	if(dot2 <= 0) return 1LL * bp.x * bp.x + 1LL * bp.y * bp.y;
	long long len2 = 1LL * ab.x * ab.x + 1LL * ab.y * ab.y;
	long long area2 = abs((long long)ab * ap);
	return (area2 * area2) / len2;
}

int main(){
	Point a{0,0}, b{4,0}, p{2,3};
	cout << "dist to segment: " << dist2PointSegment(p,a,b) << "\n";   // 9
	cout << "dist to line:    " << dist2PointLine(p,a,b) << "\n";      // 9
	return 0;
}
\end{lstlisting}
